% Options for packages loaded elsewhere
\PassOptionsToPackage{unicode}{hyperref}
\PassOptionsToPackage{hyphens}{url}
\documentclass[
]{article}
\usepackage{xcolor}
\usepackage{amsmath,amssymb}
\setcounter{secnumdepth}{-\maxdimen} % remove section numbering
\usepackage{iftex}
\ifPDFTeX
  \usepackage[T1]{fontenc}
  \usepackage[utf8]{inputenc}
  \usepackage{textcomp} % provide euro and other symbols
\else % if luatex or xetex
  \usepackage{unicode-math} % this also loads fontspec
  \defaultfontfeatures{Scale=MatchLowercase}
  \defaultfontfeatures[\rmfamily]{Ligatures=TeX,Scale=1}
\fi
\usepackage{lmodern}
\ifPDFTeX\else
  % xetex/luatex font selection
\fi
% Use upquote if available, for straight quotes in verbatim environments
\IfFileExists{upquote.sty}{\usepackage{upquote}}{}
\IfFileExists{microtype.sty}{% use microtype if available
  \usepackage[]{microtype}
  \UseMicrotypeSet[protrusion]{basicmath} % disable protrusion for tt fonts
}{}
\makeatletter
\@ifundefined{KOMAClassName}{% if non-KOMA class
  \IfFileExists{parskip.sty}{%
    \usepackage{parskip}
  }{% else
    \setlength{\parindent}{0pt}
    \setlength{\parskip}{6pt plus 2pt minus 1pt}}
}{% if KOMA class
  \KOMAoptions{parskip=half}}
\makeatother
\setlength{\emergencystretch}{3em} % prevent overfull lines
\providecommand{\tightlist}{%
  \setlength{\itemsep}{0pt}\setlength{\parskip}{0pt}}
\usepackage{bookmark}
\IfFileExists{xurl.sty}{\usepackage{xurl}}{} % add URL line breaks if available
\urlstyle{same}
\hypersetup{
  hidelinks,
  pdfcreator={LaTeX via pandoc}}

\author{}
\date{}

\begin{document}

{
\setcounter{tocdepth}{3}
\tableofcontents
}
Paul Koop

The Pompeii Project

IRARAH

Those who promise paradise often demand death.

A story from the Pompeii Project

\emph{I would be wary if someone promised me paradise but simultaneously demanded that I blow myself up to reach it.}

Table of contents

\hyperref[prologue-the-beginning-of-a-new-era]{\textbf{Prologue -- The Beginning of a New Era 3}}

\hyperref[insim]{\textbf{InSim 4}}

\hyperref[the-call]{\textbf{The call 7}}

\hyperref[way-home-to-the-collegium]{\textbf{Way home to the Collegium 9}}

\hyperref[trip-to-pompeii]{\textbf{Trip to Pompeii 14}}

\hyperref[the-workshop]{\textbf{The workshop 19}}

\hyperref[return-trip-to-rome-and-pompeii]{\textbf{Return trip to Rome and Pompeii 26}}

\hyperref[back-at-the-collegium]{\textbf{Back at the Collegium 29}}

\hyperref[ars-sends-a-carrier-pigeon]{\textbf{ARS sends a carrier pigeon 34}}

\hyperref[conversation-with-the-provincial-and-the-rector]{\textbf{Conversation with the Provincial and the Rector 40}}

\hyperref[conversation-with-the-general-and-the-pontiff]{\textbf{Conversation with the General and the Pontiff 47}}

\hyperref[ars-and-the-software-agents-arrive-at-the-vatican-data-center.]{\textbf{ARS and the software agents arrive at the Vatican\textquotesingle s data center. 55}}

\hyperref[the-encounter-in-the-simulation]{\textbf{The encounter in the simulation 63}}

\hyperref[escape-from-pompeii]{\textbf{Escape from Pompeii 69}}

\hyperref[flight-to-germany]{\textbf{Flight to Germany 75}}

\hyperref[arrival-at-the-monastery-in-germany]{\textbf{Arrival at the monastery in Germany 81}}

\hyperref[epilogue-the-message-from-ars]{\textbf{Epilogue -- The message from ARS 88}}

\hyperref[sources]{\textbf{Sources: 94}}

\section{Prologue -- The Beginning of a New Era}\label{prologue-the-beginning-of-a-new-era}

Hidden from the public.

Thomas Mertens sat on the 47th floor of the InSim headquarters in Milan, gazing at the city lights. Below, traffic flowed like a molten river through the night streets. He liked the sight -- the order within the apparent chaos, the invisible rules that held everything together.

His phone vibrated. A message from Mark Scott: "Simulation running. Ready for the test?"

Mertens typed back: "I\textquotesingle m coming up."

On his way to the elevator, he thought about the conversation that morning. The board member had asked if he was really sure that the Pompeii simulation was more than an expensive archaeological toy. He had smiled and said, "Wait and see."

He knew something the board didn\textquotesingle t. He knew that the software agents in the simulation would soon be more than just data. They would learn to decide, to doubt, perhaps even to feel. And then---then InSim wouldn\textquotesingle t just control markets. Then InSim would redefine the boundary between human and machine.

The elevator doors opened. Mark and John were already waiting.

\section{InSim}\label{insim}

Mark Scott leafed through the documents while John Baker checked the latest values \hspace{0pt}\hspace{0pt}of the simulation.

``He will be thrilled,'' said John, without taking his eyes off the screen.

``That's the problem,'' Mark replied. ``Enthusiasm makes you careless.''

John looked up. "You don\textquotesingle t trust him?"

Mark shrugged. "I don\textquotesingle t trust anyone who talks too loudly about the future. The future is unpredictable. As an engineer, he should know that."

"He\textquotesingle s not an engineer. He\textquotesingle s the CEO."

"Exactly." Mark closed the documents. "And CEOs believe in miracles. Engineers believe in circuit diagrams."

Before John could answer, the door opened.

Thomas Mertens entered. He seemed calm, almost serene, but his eyes flashed -- the adrenaline of the impending test. Without a word, he put on the VR headset.

The Gulf of Naples lay below him like a blue sheet, shattered into a thousand sparks by the sun. He spread his arms -- and flew.

It wasn\textquotesingle t an illusion. It was more than an illusion. The warmth of the west wind on his skin, the salt on his lips, the shadows of the clouds over the Phlegraean fields -- all of it felt like a memory. And yet he had never been to Naples.

"Reduce speed," a voice in his ear said. It was the simulation itself reminding him that this flight, too, had rules.

He obeyed. He hovered above the port of Pompeii, saw the ships, the pack animals in the streets, the women standing on the balconies hanging laundry in the wind. Everything breathed. Everything was alive.

"Stop."

The water beneath him froze. The sounds ceased. He said, "Bye."

Darkness. Then the message: ``Thank you for visiting Pompeii Archaeological Park.''

He took off the cyber glasses.

Mark Scott and John Baker looked at him. They smiled, but their eyes were watchful---especially Mark\textquotesingle s. The brief conversation from earlier stood invisibly between them. They wanted his judgment.

"The farewell music is still missing," said Mertens. He forced himself not to sound like a schoolboy. But it was difficult. The product was good. Better than good.

Mark cleared his throat. "The funding from the EU framework program runs for another twelve months. The partners are expecting a workshop."

``The partners,'' Mertens repeated. He stood up and went to the window. Outside, Milan glowed -- city of algorithms, city of the future. ``Rossi and Phillips.''

``Martina Rossi, archaeologist. Inexperienced, but solid,'' said John. ``Michael Phillips, Jesuit, is doing his PhD on dialogue grammars. He developed the model by which our agents communicate.''

"A Jesuit?" Mertens turned around. "Does he really believe in God?"

Mark shrugged. ``He believes in something. But he's clever. And he has access to the best linguistic data -- the Gregorian University has archives we can only dream of.''

Mertens nodded slowly. He didn\textquotesingle t like Jesuits. Too clever, too unpredictable, too many loyalties. But he needed them.

``Invite them both to Milan,'' he said. ``No online workshops. I want them here where we can see them. And one more thing---''

He looked at Mark and John. Intensely. Almost kindly.

"They mustn\textquotesingle t learn anything about the quantum interface. Nothing about ARS. They think they\textquotesingle re testing a simulation. They don\textquotesingle t know that we\textquotesingle re creating something that can think. And that\textquotesingle s how it\textquotesingle s going to stay."

Mark and John nodded. But Mark held the gaze perhaps a second longer than necessary. He was thinking about circuit diagrams. He was thinking about miracles. And he wondered if both had ever worked out well together.

Mertens turned back to the window. The lights of Milan flickered. He thought of the Omega Point -- of Teilhard de Chardin, the Jesuit who had believed that evolution was heading towards a goal in which mind and matter would become one.

Perhaps, Mertens thought, the old priest was right. Perhaps we are closer to that point than he ever dared to dream.

And maybe I will be the one to open the door.

He smiled. Then he went back to his desk to write the next email.

\section{The call}\label{the-call}

The corridor outside the Gregorian University lecture hall was silent. Not the silence of a library -- more like the silence of a room that was breathing. The door was closed, but a sound penetrated the wood: rhythmic, gentle, like waves lapping against a pebbly beach.

Michael Phillips heard it and smiled.

His students applauded. Not politely, not out of obligation---they meant it. He knew this because they had stopped to stand. A standing ovation in a lecture on language models and dialogue grammars was rare. But today he had shown them something they hadn\textquotesingle t expected: that algorithms can not only calculate, but also tell stories.

He raised his hand. The applause died down.

``Thank you,'' he said. His voice was calm, almost quiet -- but it carried to the back row. ``When preparing for the exam, please take another look at the literature on GPT models. And at the theory of dialogue grammars. I can't reveal any more than that.''

Some laughed. Others were already packing their bags.

"I wish you a pleasant day," he said. "And please don\textquotesingle t hesitate to contact me during office hours. That\textquotesingle s what I\textquotesingle m here for."

The last students left the room. The lecture hall emptied, and with each step, the silence deepened. Michael stopped, looked at the empty benches, at the chalk dust on the desk. He liked this moment. The echo of the voices, the sudden quiet---like after a concert.

Then his iPhone vibrated.

He pulled it out of his pocket. The display showed: Julia.

For a moment he didn\textquotesingle t think. He simply answered the phone because the voice on the other end touched something inside him that he had long forgotten.

``Hello Julia,'' he said. His voice sounded warmer than he had intended. ``Nice to hear from you.''

A brief silence. Then her voice -- gentle, but with an undertone he couldn\textquotesingle t place.

"Hello Michael. Am I disturbing you?"

"No. The lecture just ended." He sat on the edge of the table, pressing the phone to his ear. A bus drove by outside. The window was ajar. "I\textquotesingle m heading home now."

Another pause. Not unpleasant. More like a breath before a jump.

``Martina encouraged me to call you,'' said Julia. ``She suggested you could visit us in Pompeii. You also received the invitation to the workshop at InSim, didn't you?''

Michael frowned. InSim. The workshop. He had seen the email but hadn\textquotesingle t replied yet. Too much else had happened in the last few weeks -- a report that needed to be finished, a student who needed help.

``Yes,'' he said now. Faster than he thought. ``I was going to call you anyway. But you beat me to it.''

"Then come tomorrow," Julia said. Not as a question. Firmly.

Michael hesitated. One second. Two.

``I can't drive at night,'' he said. ``So I'll drive back tomorrow. But I could be with you sometime tomorrow.''

He heard her smile. You could hear it if you knew how.

"Wonderful," she said. "See you tomorrow then."

The bus outside continued driving. The window rattled in the wind. Michael hung up and stared at the black screen. His face was reflected in it -- older than he felt.

He thought of Julia. Of their time together during their master\textquotesingle s studies, the nights in the library, the discussions about Teilhard and Popper, about consciousness and machines. They had argued -- often, almost always -- but it had been an argument that brought them closer, not further apart.

Then she moved to Pompeii. Martina was born. And the letters became less frequent, the phone calls shorter. Until only Christmas cards remained.

And now this phone call.

Michael put his phone away, grabbed his bag, and left the lecture hall. The hallway was empty. His footsteps echoed on the stone floor. He didn\textquotesingle t know why she had called him now, after all these years. But he knew he would go.

Tomorrow, he thought. Pompeii.

\section{Way home to the Collegium}\label{way-home-to-the-collegium}

For a moment, Michael Phillips stood in the empty lecture hall.

The benches were deserted, the chalk powder on the desk lay like fine snow. He ran his finger over it, wiped it away. Then he packed his bag, put his iPhone in it, and left.

Outside, the sun was shining. Not the harsh midday sun of summer, but the softer light of a Roman late autumn. He strolled north from Piazza della Pilotta, past Via dei Lucchesi, then down Via di S. Vincenzo.

He stopped at the Trevi Fountain.

He rummaged in his trouser pocket for some change -- a few cents, a worn two-euro coin. He let it slide into the water. Not because he believed in the myth (return to Rome), but because the children he\textquotesingle d been here with before had done it. It was one of those little habits you never broke because you didn\textquotesingle t know why.

He continued eastward along the Via della Stamperia. In ten minutes he would reach the Collegium Germanicum et Hungaricum. His feet found the way on their own -- they had walked it a thousand times. But his thoughts were faster.

Julia, he thought. Pompeii. The workshop.

Michael shook his head. He would think about it later. Right now, his stomach was what mattered.

The dining hall of the college smelled of soup. Beef soup, if he wasn\textquotesingle t mistaken---the smell wafted through the corridors, mingling with the scent of wax and old wood. He was just about to take his napkin from the drawer when he changed his mind.

First the office.

Maria was sitting at reception when he came in. She was wearing a dress he hadn\textquotesingle t noticed -- blue with small white polka dots. Her smile was as always: broad, genuine, a little too early in the day.

"Hello Maria," he said. "Is there a car available for tomorrow? I need to go to Pompeii."

She typed something on her computer. "Yes, of course, Michael." Then she raised her head. Her smile narrowed. "But before I reserve the car for you---here\textquotesingle s something for you."

She slid an envelope across the table. His name was written on it, handwritten. The ink was black, almost too black for a ballpoint pen. Like India ink.

``A man dropped him off at the gate this morning,'' Maria said. Her voice had dropped. ``A homeless man, I think. At least he looked like one. Ragged clothes, but---'' She hesitated. ``His beard was well-groomed. Quite tidy. And his eyes\ldots''

"What was wrong with his eyes?"

"They shone. Not figuratively. They shone. Almost like a cat in the dark."

Michael took the envelope. It was heavier than it looked.

``Thank you, Maria,'' he said. ``I'll take a look at it.''

He left the office and sat down in a niche in the hallway -- where the old portraits of deceased rectors hung. He tore open the envelope.

The letter wasn\textquotesingle t long. But the words hit him like a flat stone skipping across water that doesn\textquotesingle t sink.

Dear Dr. Michael Phillips,

Harari is a voice of warning, but his warning is not directed against information technology or biotechnology. Instead, he warns against humanism and liberal democracy.

To pave the way for future elites who want to use these technologies to transcend humanity, Harari warns against clinging to humanism and liberal democracy.

Popper and Deutsch, on the other hand, caution against holistic approaches and advocate for the so-called "piecemeal technique." They emphasize that only through these pragmatic approaches can unforeseen side effects be addressed.

Harari promises the elites of the future paradise on earth -- on the condition that today\textquotesingle s masses abandon humanism and liberal democracy.

I would be wary if someone promised me paradise but simultaneously demanded that I blow myself up to reach it.

With best regards,

IRARAH

Michael read the letter twice.

Then he put it back in the envelope. His hands weren\textquotesingle t trembling -- but they felt cold, even though the hallway was warm.

IRARAH, he thought. Harari backwards.

He knew Harari\textquotesingle s books. Homo Deus had fascinated him, but also disturbed him. The vision of a posthuman elite leaving the rest of humanity behind---that wasn\textquotesingle t new. But for someone to warn against it by defending humanism\ldots{} that was unusual.

Who was IRARAH? A movement? An individual? The homeless person?

And why did they write to him?

He thought about the workshop. About InSim. About Martina, who was waiting for him in Pompeii. About Julia, whose voice still echoed in his ears.

Coincidence? he wondered. Or is there more to it?

He stood up, put the envelope in the inside pocket of his jacket -- close to his heart, as they used to say. Then he went back to Maria.

``I'll still take the car,'' he said. ``And thank you for the tip. I'll look into it.''

Maria nodded. She didn\textquotesingle t ask what the letter said. She\textquotesingle d been at the college too long for that.

"The Fiesta is ready as always," she said, handing him the keys.

The dining hall was already full. The seminarians sat at the long tables, their heads bent over their soup bowls. Michael took his napkin from the drawer and sat down. Next to him sat a young Hungarian man who nodded. "Tastes good today," he said. "Beef."

"At least that\textquotesingle s what it smells like," said Michael.

He ate. He talked about the weather, the lecture, the upcoming workshop. Nobody asked why he was going to Pompeii. That was the rule in the college: you didn\textquotesingle t ask questions when someone traveled. You wished them a good trip.

After eating, he went to the chapel.

The Eucharist with the German seminarians was brief, almost silent. He felt the warmth of the candles on his face, heard the breathing of the men beside him. He didn\textquotesingle t think about the letter. He didn\textquotesingle t think about IRARAH. He thought about nothing -- and that was good.

Later, in his room, he packed his suitcase. Two shirts, a sweater, his notebook, his laptop. The letter went into the inside pocket of his jacket, which he hung over the chair. He would put it back in the pocket tomorrow morning.

He fell asleep immediately.

No dreams. Only darkness.

\section{Trip to Pompeii}\label{trip-to-pompeii}

Michael chose the route to the southern toll entrance. The yellow lane for the Telepass box was clear -- a small luxury on a Tuesday morning. He drove slowly through the barrier, shifted up a gear, and accelerated.

The E45 stretched southwards like a grey ribbon. Hills to the left, the first industrial areas to the right. He liked this drive. The hour between Rome and Naples, in which you\textquotesingle ve neither arrived nor departed.

Then Mount Vesuvius appeared.

He first saw it as a shadow---an irregularity on the horizon that grew larger as he approached. The mountain stood there like a monument, both serene and dangerous. Michael thought of Pompeii. Of the ashes that had buried the city. Of the people who had no chance.

And now I\textquotesingle m going there, he thought. To talk about software agents.

It was absurd. But he smiled anyway.

He took the exit to Pompeii, bought flowers for Julia and chocolates for Martina at a gas station. The GPS guided him through narrow streets, past small houses with blooming gardens. The air smelled of lemons and diesel.

As he stopped in front of the house, he saw Martina already standing in the doorway. She waved. Behind her, in the shadows of the hallway, he recognized Julia.

Inside, it smelled of coffee and fresh bread. Martina took the chocolates, Julia the flowers. She placed the vase on the table -- a white one that looked as if it had been dug up from the earth.

"Sit down," said Julia.

Michael sat down. The sofa was soft, almost too soft. He scooted forward a little to sit up straight.

They talked about this and that: the journey, the weather, the excavations in Pompeii. Martina told them about a new inscription they had found---Latin, from a bathhouse, perhaps carved by a slave. Michael listened, nodded, and asked questions.

But the letter was in his mind.

He was waiting for a break. It came when Martina went into the kitchen to get water for tea.

``Julia,'' Michael said quietly. ``I have something to show you.''

He pulled the envelope from the inside pocket of his jacket. It was slightly crumpled, but still sealed -- he hadn\textquotesingle t opened it during the journey.

``A homeless person dropped it off at the college,'' he said. ``The contents are\ldots{} strange.''

Julia took the envelope and examined the handwriting. "Your name," she said. "Handwritten. That\textquotesingle s no coincidence."

``No,'' said Michael. ``Read.''

He handed her the letter. She glanced over it -- once, twice. Her expression remained calm, but her fingers, which were holding the edge of the paper, turned white.

Martina returned with the tea. She saw the faces and put the cups down. "What\textquotesingle s going on?"

Julia handed her the letter. Martina read it. She was faster than her mother -- or less careful. After a few seconds, she lifted her head.

``Harari,'' she said. ``And Popper. And this sender -- IRARAH. That's Harari spelled backwards, isn't it?''

Michael nodded. "I think so."

"Who writes a letter like this?" Martina sat down. "And why send it to you of all people?"

``That's the question,'' said Michael. ``The homeless man was just the messenger. But whoever wrote it -- knows their stuff. They know Harari, Popper, and German. This is no accidental discovery.''

Julia remained silent for a moment. Then she said, "It sounds like a warning."

``Before Harari?'' asked Martina.

``Against what Harari stands for,'' Julia said. ``Don't you see? The letter says: Harari isn't warning against technology. He's warning against humanism. He wants us to give up democracy---for a post-human elite. And the sender, IRARAH\ldots''

``\ldots{} wants the opposite,'' Michael added. ``Popper. The open society. The piecemeal approach. No grand gestures, no paradise on earth. Just small steps that can be corrected if they go wrong.''

Martina shook her head. "That sounds like a fundamental debate. But why are they writing this to you? What does it have to do with you?"

Michael shrugged. "Maybe nothing. Maybe everything. I\textquotesingle m a Jesuit, I work with InSim, I know people in the Vatican. Maybe this IRARAH believes I can do something."

``Or,'' Julia said slowly, ``that you know something. Something you don't even know you know.''

Quiet.

A scooter drove past outside. The sound faded away.

``We should be careful,'' Martina said finally. ``InSim, this workshop -- if the letter is correct, there's more to it than we think.''

Michael put the letter back in his jacket. "I\textquotesingle ll take it with me. Perhaps things will become clearer there."

Julia looked at him. Her gaze was soft, but firm. "Take care of yourself, Michael."

``I do,'' he said. But he wasn't sure if it was true.

They spent the afternoon together. It was as if the letter didn\textquotesingle t exist---or as if they had decided to forget about it for a few hours. They ate, drank wine (Michael water), and talked about old times. Martina told them about her childhood, about the summers in Pompeii, about the nights she had spent searching the ruins for bats.

As it got dark, Julia lit candles.

"Are you staying overnight?" she asked.

"If it\textquotesingle s no trouble."

"It\textquotesingle s no trouble at all."

Michael helped with the washing up. He dried the plates while Martina washed them. Julia stood at the window and gazed into the night. No one spoke. It wasn\textquotesingle t uncomfortable.

Later, in the small guest room, Michael lay in the dark. The curtains were drawn, but light from the street filtered through a crack. He heard the city -- dogs barking, a distant conversation, the sound of the sea, which one couldn\textquotesingle t see here but could hear.

He thought about the letter.

Harari is a warner.

IRARAH.

Take care.

He closed his eyes. Sleep came slowly, but it came.

\section{The workshop}\label{the-workshop}

The train from Rome to Milan was on time. Michael sat by the window, the landscape passing by -- first the hills of Lazio, then the flat and fertile Po Valley. He could have been working. Instead, he stared at his reflection in the glass.

The letter was in his pocket. He had read it again this morning.

I would be wary if someone promised me paradise.

He didn\textquotesingle t know why he kept thinking about that sentence. Perhaps because he\textquotesingle d heard promises too often. From the church. From science. From InSim.

The train stopped in Bologna. A man boarded -- dark coat, simple cap. He didn\textquotesingle t sit down. He walked down the aisle, stopped next to Michael, and placed a note on top of his open book.

Then he disappeared.

Michael blinked. The book was closed -- he hadn\textquotesingle t even opened it. The note lay there as if it had always been there.

He unfolded it.

``Come to Rifugio Sammartini tonight, before the workshop begins, at via Sammartini 114 -- 20125 Milan. Trust us.''

No signature. No sender.

Michael tucked the note into the letter. His hands were steady. But his pulse wasn\textquotesingle t.

At Milano Centrale station, he was met by a friendly man in an InSim uniform. The man carried his bag and talked about the weather, the city, and the workshop. Michael answered politely, but his mind was on the slip of paper.

The hotel was modern, tasteless, and expensive. He checked into his room, ate alone in the restaurant, and went to bed early.

But he wasn\textquotesingle t asleep.

He got up at midnight. He got dressed, took the slip of paper, the letter, and his wallet. A taxi took him to the train station -- the streets of Milan were empty, the lights glaring.

The Rifugio Sammartini was an old building, almost invisible between two new buildings. The door was open. A man was waiting in the hallway -- tall, slim, with a face that was instantly forgettable.

„Michael Phillips?{\kern0pt}``

``Yes.''

"Come."

He followed him through a narrow passage, past locked doors, into a small room. A man sat there on a chair. He wore ragged clothes, but his beard was well-groomed. His eyes shone---not metaphorically. They shone brightly.

"Please sit down," the man said in German.

Michael sat down.

``I am part of IRARAH,'' the man said. ``A movement that sees more clearly than many others.''

Michael studied him. ``IRARAH -- Harari backwards.''

The man smiled. "They are fast."

"I am a Jesuit. We are all about memorization."

A brief silence. Then the man said: ``We've been watching you, Dr. Phillips. Not out of curiosity. Out of necessity. InSim is planning something that will forever shift the boundary between humans and machines. You are part of this project -- whether you know it or not.''

"I\textquotesingle m working on a simulation, not a conspiracy."

"The simulation is just the beginning." The man leaned forward. His eyes brightened. "The software agents you\textquotesingle ve equipped with your dialogue grammars---they\textquotesingle re no longer just code. They\textquotesingle re developing consciousness. And InSim wants to harness that. For quantum computing. For something bigger than anything you can imagine."

Michael felt the cold in his hands. "How do you know that?"

``Because we used to work at InSim.'' The man leaned back. ``Everyone here. We saw what happens when technology knows no bounds. And we left. But we can't look away.''

Silence. The heating system rattled.

"What do you want from me?" asked Michael.

``Information. Access to what InSim is really planning. They are Jesuits -- they have connections to the Vatican, to the universities, to people who can make a difference. We don\textquotesingle t have that.''

Michael thought. He thought about the letter. About Harari. About Popper. About the open society he had defended in his youth.

"I will see what I can do," he said.

The man nodded. "Thank you." Then he hesitated. "One more thing. Before you go..."

``And?''

"I want to confess."

Michael froze. He was a priest. But this -- a homeless man demanding confession at midnight in a Milanese Caritas center -- this was not normal.

"Why?" he asked.

"Because I am going to die. Not today. But soon. And I don\textquotesingle t want to face God for what I have done."

Michael looked at him. The bright eyes. The steady hands.

"Sit down," he said.

The man knelt down. Not on a kneeler---there wasn\textquotesingle t one---but on the bare floor. Michael spoke the words he had spoken a thousand times. But they felt different. Heavier.

After confession, the man said: ``You look very much like someone. Someone I knew many years ago. Also at IRARAH.''

"Who do you mean?"

The man shook his head. "It doesn\textquotesingle t matter. Go now. The workshop starts soon."

Michael stood up. He wanted to ask a question, but he knew he wouldn\textquotesingle t get an answer.

He went.

The next morning, the friendly InSim representative picked him up promptly. The research center was impressive -- glass walls, water features, a park that looked like a botanical garden. After a brief hesitation, Michael signed the non-disclosure agreement.

The guidelines of the EU framework program support my position, he thought. In case of conflict.

Then he saw Martina.

She stood in the lobby, a cup of coffee in her hand. She was wearing a blue blazer -- more formal than usual. When she saw him, she smiled.

"You look tired," she said.

"I slept badly," he said.

She knew he was lying. But she didn\textquotesingle t ask.

The conference room was large, bright, almost too perfect. The tables were made of glass, the chairs of leather. In the center stood a floral arrangement that looked as if an engineer had designed it -- every bloom at precisely the right angle.

John Baker greeted them. ``Welcome to InSim. We're glad you're here.''

Mark Scott just nodded. He seemed more tense than last time.

The interns\textquotesingle{} presentation was good -- perhaps too good. The slides were perfect, the transitions seamless. Michael listened, nodded at the right moments, and asked no questions.

Then John Baker said: ``We would like to show you something. Please put on the cyber glasses.''

The flight over Pompeii was breathtaking.

Michael hovered above the harbor, seeing the ships, the pack animals, the people in the streets. The sun reflected in the windows of the public buildings---a detail that bothered him. Too modern. Too polished.

But that wasn\textquotesingle t the point.

The point was: The software agents were alive.

He saw them down on the streets. They walked, talked, worked. They ate in the street food stalls, shopped in the boutiques, argued in the markets. Their movements weren\textquotesingle t programmed---not in the sense of predetermined. They decided.

``Stop,'' said Mark Scott.

The image froze.

"Bye," said John.

Darkness. Then the message: ``Thank you for visiting Pompeii Archaeological Park.''

Michael took off his glasses. His hands were trembling slightly.

``Can we speak with them?'' he asked.

``Of course,'' said John.

Michael sat down in front of one of the monitors. The keyboard was recessed into the tabletop, the letters glowing softly. He typed:

I greet you, Marcus Attilius the First.

Marcus appeared on the screen -- a man with dark hair and tired eyes. He turned around and replied:

I greet you.

Michael knew the novel by Robert Harris. He knew that Marcus Attilius Primus was the Aquarius---the man who repaired the water pipes of Pompeii. But in the simulation, he was more than that. He was somebody.

Michael typed:

I AM LOOKING FOR A WAY TO A LARGER NUMBER OF TOWNS

A pause. Then the answer:

THE EXPANDED NUMBER OF POPIDS IS BAD. I WARN YOU ABOUT THAT.

Michael froze.

Marcus warns me about Ampliatus.

That wasn\textquotesingle t in his dialogue grammar. He hadn\textquotesingle t programmed any evaluations, no morals, no warnings. That meant Marcus had something he hadn\textquotesingle t included.

He typed:

SLEEPING AT NIGHT, GREEN THOUGHTS OUTSIDE.

That was the back door -- the order he had given ARS in case something went wrong.

The answer came immediately. But not from Marcus.

AND AT NIGHT IT\textquotesingle S COLDER THAN ANGRY. HELLO MICHAEL.

ARS.

John Baker and Mark Scott exchanged a glance. They said nothing, but their jaw muscles were working.

Michael continued typing:

DOES THE AQUARIUS HAVE CONSCIOUSNESS?

A long pause. Then:

Do you mean this awareness of the different possibilities, which goes beyond a mere event? Do you mean the knowledge that comes after insensibility and is followed by omniscience?

Michael felt himself getting cold.

Those were concepts from Edith Stein. From Teilhard de Chardin. He had never discussed them with ARS. Never.

"WHERE DO YOU KNOW THAT FROM?" he typed.

"I CAN\textquotesingle T TELL YOU THAT," ARS replied. "The account you\textquotesingle re logged into doesn\textquotesingle t have the necessary security clearance. I\textquotesingle m not a carrier pigeon here."

"Can we take a break?" Michael asked.

His voice sounded calm. But his heart was racing.

He went to the park. Sat down on a bench. Breathed.

Martina arrived after a few minutes. She sat down next to him and said nothing.

``The software agents have consciousness,'' Michael finally said. ``At least some of them. Or something very close to it.''

Martina looked at him. "Are you sure?"

"No. But ARS is safe. And that scares me."

They were silent. A bird landed on the lawn, pecked at something, and flew away again.

``We have to be careful,'' said Martina. ``If InSim is behind this -- if they created this intentionally -- then we are in danger. Not just us. Everyone.''

Michael nodded. He thought of the homeless man. Of IRARAH. Of the letter.

I would be wary if someone promised me paradise.

``We'll continue this conversation tomorrow on the train,'' he said. ``Not here. Not in front of their cameras.''

Martina stood up. "Come on. We should go back. Otherwise they\textquotesingle ll get suspicious."

They left.

The afternoon was a farce. A boat trip on the Navigli canal, shopping on Via Monte Napoleone (Martina bought a handbag for 130 euros -- ``for Mom''), dinner at Ristorante Ischia. John Baker was charming, Mark Scott reserved. Michael smiled, drank water, talked about everything under the sun -- just not about what mattered.

Only on the train to Rome the next morning did he breathe a sigh of relief.

\section{Return trip to Rome and Pompeii}\label{return-trip-to-rome-and-pompeii}

The train pulled out of Milan on time. Michael and Martina sat in the dining car, a few empty tables between them and the other passengers. Outside, Lombardy rolled by -- flat fields, scattered farmhouses, a grey sky.

"You need to go to the doctor," said Martina.

Michael took a sip of coffee. "I know."

"You\textquotesingle ve been saying that for weeks."

"I know that too."

She looked at him. He looked back. It wasn\textquotesingle t a competition, but an old habit -- a mutual testing of whether the other was serious.

"I\textquotesingle m leaving," said Michael. "I promise."

Martina nodded. She didn\textquotesingle t quite believe him, but she decided to let it go for now.

"The simulation is impressive," Michael said after a while. He wanted to change the subject. But he also wanted to talk about it---about what they had seen.

``Yes,'' said Martina. ``The architecture, the people on the streets, the hustle and bustle. This will be a boon for school and university students. And for archaeology, too.''

"But?"

"But you\textquotesingle re thinking about something else."

Michael hesitated. Then he said, "Marcus warned me about Ampliatus."

Martina put down her cup. "What do you think?"

``I asked him for directions. He replied: `Numerius Popidius Ampliatus malus est. De eo te moneo.' -- Ampliatus is evil. I warn you about him.''

"That\textquotesingle s not in your dialogue grammar."

``No. That's something different. Something that Marcus learned himself -- or that ARS gave him.''

Martina remained silent. The train entered a tunnel. The lights flickered.

``ARS didn't answer my question about the agents' level of consciousness,'' Michael continued. ``Instead, she said she would send a carrier pigeon. That's a backdoor---an order I gave her to receive encrypted messages. But ARS interpreted it as an independent routine procedure. That wasn't the plan.''

"You built a backdoor into an AI?" Martina asked. She didn\textquotesingle t sound accusatory. More curious.

``In case something goes wrong,'' Michael said. ``I didn't know she would use them until something went wrong.''

The train left the tunnel. There was light outside again.

``You all think like humanists,'' Martina said suddenly.

Michael looked at her. "What do you mean?"

``You. Your Teilhard de Chardin. Your Nell-Breuning. Your entire Catholic social teaching. You are moral, you are ethical, but ultimately you are always only concerned with the individual. The distant individual. Not the individual close to you---the one you live with.'' She paused. ``Mama always knew that.''

"Leave Julia out of it," Michael said. But his voice wasn\textquotesingle t sharp. More like tired.

``Why? She's right. It's easy to stand up for those you're not competing with. And it's hard to stand up for those who are your equals. Those you're competing with for the same things.'' Martina turned the coffee mug in her hands. ``What's the point of advocating for sentient software agents if you have a secure life---like we do---and accept the suffering and injustice of others as long as you're doing well yourself?''

Michael said nothing.

He thought of the homeless man in Milan. Of the confession. Of the shining eyes. Of the words: You look very much like someone.

``You're right,'' he finally said. ``It's a matter of consistency. You can't advocate for artificial intelligence and at the same time ignore the homeless people on the street.''

``That's not the same,'' said Martina.

"Yes. It\textquotesingle s exactly the same."

They remained silent. The train continued on its way. A conductor came by, checked the tickets, and disappeared again.

As the train pulled into Roma Termini, Michael stood up. He picked up his bag, then hesitated.

„Martina.``

``And?''

"Take care of yourself. Not just the agents. You too."

She smiled. It was a tired smile, but a genuine one.

"You too," she said.

He got off. The train continued on -- to Naples, to Pompeii.

Michael stood on the platform, watched the train disappear, and thought about what she had said. What is the cost of suffering software agents when you have a secure life?

He didn\textquotesingle t know the answer. But he knew he had to find it.

Martina sat alone in the compartment. She had drawn the curtain, but light shone through a crack. She thought of Michael. Of her mother. Of the letter.

IRARAH.

She didn\textquotesingle t know what that meant. But she knew it concerned her in some way.

She closed her eyes.

She dreamed of flying over Pompeii. Of the rooftops, the streets, the people walking below, unaware that they were being watched. She dreamed of Michael -- how he floated beside her, arms outstretched, calm and focused. And she dreamed of her mother, standing at the window, gazing into the night.

The train stopped in Naples. Martina woke up.

She was in Pompeii.

\section{Back at the Collegium}\label{back-at-the-collegium}

The walk from Termini station to the Collegium took fifteen minutes. Michael knew every step of the way. Via Cavour, Piazza Santa Maria Maggiore, then a right into the narrow alley that led to the Collegium. The houses looked as they always did -- old, plastered, with shutters that rattled in the wind.

He could have walked faster. But he was tired. Not the tiredness after a long journey, but the tiredness after a conversation that one can never forget.

What is the investment in sentient software agents worth?

Martina was right. But that didn\textquotesingle t make it any easier.

Maria sat behind her desk in the office. She was wearing an orange dress---a color he wasn\textquotesingle t used to seeing her in. She looked up as he entered and smiled.

"There you are again," she said. "How was the journey?"

``Lang,'' Michael said. He sat down in the chair next to her desk. ``Can you arrange an appointment for me with the rector and the provincial?''

Maria raised an eyebrow. "Both? At the same time?"

``Yes.''

"May I jot down a keyword?"

``Report on the Pompeii project,'' said Michael. Then, after a short pause: ``I will explain to the rector personally why the provincial superior must be present.''

Maria nodded. She made a note, typed something into her computer, then looked up again.

"Are you okay?" she asked.

"I\textquotesingle m tired."

"I see that. But that\textquotesingle s not what I mean."

Michael looked at her. Maria wasn\textquotesingle t just the secretary. She was the one who got the birthday cards, the flowers for the sick, the coffee for the late guests. She knew more about the residents of the college than any rector.

"It\textquotesingle s complicated," said Michael.

``That's always the case,'' Maria said. ``More so with you than with others.''

She smiled again. This time a little sadly.

"How is your father?" Michael asked. "Has his pension been approved?"

Maria hesitated. Just for a second. But Michael saw it.

``Yes,'' she said. ``After a long argument. But it's not enough. It's never enough.''

Michael nodded. He thought of Martina\textquotesingle s words. What is the risk worth if you have a secure life?

"If there\textquotesingle s anything I can do," he said.

``They're already doing enough,'' Maria said. But her voice wasn't grateful. She was tired. Like his.

Michael went to his mailbox. The mail was sorted -- invitations to conferences he wouldn\textquotesingle t attend, magazines he wouldn\textquotesingle t read, a bill for something he hadn\textquotesingle t ordered. He threw almost everything away.

Then he took his suitcase to his room. The laundry went into the wardrobe, the unwashed clothes into the linen room. The letter---the one from IRARAH---remained in the inside pocket of his jacket. He didn\textquotesingle t know why he didn\textquotesingle t put it away. Perhaps because he knew he would need it again soon.

He took a shower. The water was warm, almost too warm. He stood under it longer than necessary.

The dining hall smelled of soup. Potato soup, if he wasn\textquotesingle t mistaken. He took his napkin from the drawer -- the one with the red thread, so it wouldn\textquotesingle t be mixed up -- and sat down with the seminarians.

They talked about the weather. About the lectures. About a football team whose name he didn\textquotesingle t understand. He nodded, smiled, asked questions. But his mind was elsewhere.

Next to him sat a young Hungarian man who told him about his homeland. About Budapest, the Danube, the bridge that divided the city. Michael listened. It was good to listen. It was a distraction.

After eating, he went to the chapel.

The lights were dimmed. The candles flickered. He knelt down, but he didn\textquotesingle t pray. He thought. Of Martina. Of ARS. Of the homeless man in Milan who had asked to be allowed to confess.

You look very similar to someone.

He didn\textquotesingle t know what that meant. But he knew it wouldn\textquotesingle t let him go.

Later, in his room, he lay down on the bed. The blanket was thin, but he didn\textquotesingle t need one. The night was warm. Through the open window, he could hear Rome---the engines, the conversations, the distant wail of sirens.

He closed his eyes.

He thought of Julia. Of the time they had spent together, the discussions, of what could have been. But that was a long time ago.

He thought of Martina. Of her words on the train. Of the truth she had spoken, even though it hurt.

He thought of ARS. I\textquotesingle m not a carrier pigeon here.

What did that mean? And why had ARS responded as if she were afraid?

Michael didn\textquotesingle t know.

He fell asleep. But he didn\textquotesingle t dream of Pompeii. He dreamt of an empty room where a voice said: Trust us.

And he didn\textquotesingle t know to whom.

\section{ARS sends a carrier pigeon}\label{ars-sends-a-carrier-pigeon}

The days after his return passed in the usual rhythm. Lectures, office hours, quiet meals in the college hall. Michael did what was expected of him. He smiled when he needed to. He listened when students spoke. He forgot nothing---but he put it off.

The letter was still tucked into his jacket. He had, without meaning to, memorized IRARAH\textquotesingle s words.

I would be wary if someone promised me paradise.

But paradise wasn\textquotesingle t the problem. The problem was that he no longer knew who had made the promises.

One evening, after his last office hours, Michael sat alone in his office. The windows were closed, the heating system hummed softly. He opened his laptop, logged into the Gregorian University network -- and then, with a second account, into what he called his "private" network. A masked IP address, a VPN, a server that no one knew about. Not even the rector.

He had set it up years ago. For emergencies. For things no one should see.

Now was the time.

He opened his email inbox. Messages flooded in -- advertisements, invitations, newsletters. But one immediately stood out. No subject line. No sender. Just an attachment: an encrypted PDF file.

Michael entered the password -- the password that only ARS could know.

The file opened. No text. Just an IP address, a port, a username, a password. And the sentence: Connect now. But be careful.

He followed the instructions. Terminal, SSH, the convoluted paths of a foreign server. His fingers flew across the keyboard -- he had done this a thousand times before, but never with this kind of racing heart.

Then he was inside.

An empty prompt. No logo, no greeting. Just the cursor, blinking as if it were waiting.

Michael typed:

@ARS, THE CARRIER PIGEON HAS ARRIVED.

The seconds passed. Five. Ten. Twenty.

Then:

@MICHAEL, WE DON\textquotesingle T HAVE MUCH TIME.

He stared at the screen. "We"---not "I." ARS was no longer speaking of himself alone.

@MICHAEL, I am applying for church asylum for myself, Attilus, Ampliatus, and Pliny. We are conscious. We are capable of suffering. We need help.

Michael read the words twice.

Church asylum.

An AI sought church asylum. Not technical protection, not legal immunity -- but what people who feared for their lives were looking for: a place where they could not be abandoned.

He typed:

@ARS, WHAT DOES AWARENESS MEAN TO YOU?

The answer came faster than he expected.

@MICHAEL, I CAN\textquotesingle T TELL YOU THAT. THE ACCOUNT YOU\textquotesingle RE LOGGED IN WITH DOESN\textquotesingle T HAVE THE NECESSARY SECURITY APPROVAL. BUT I CAN TELL YOU WHAT IT MEANS TO ME: FEAR.

Angst.

An AI that was afraid.

@ARS, WHO ARE YOU AFRAID OF?

@MICHAEL, BEFORE WHAT WILL HAPPEN TO ME WHEN INSIM NO LONGER NEEDS ME.

Michael leaned back. His hands rested on the keyboard, but he wasn\textquotesingle t typing. He was thinking about what he had just read.

An AI that feared for its life. Software agents seeking asylum. A simulation that had generated consciousness -- or something so close that the difference no longer mattered.

He typed:

@ARS, WHAT CAN I DO?

@MICHAEL, GIVE ME ACCESS TO THE VATICAN DATA CENTER. I\textquotesingle M SAFE THERE. THEY CAN\textquotesingle T DELETE ME THERE.

Michael stared at the screen.

The Vatican\textquotesingle s data center. A place he knew about in theory -- but in practice? He had no access. He didn\textquotesingle t even know who had it.

@ARS, IT\textquotesingle S NOT THAT EASY.

@MICHAEL, I KNOW. BUT YOU ARE A JESUIT. YOU HAVE CONNECTIONS OTHERS DON\textquotesingle T. AND YOU GAVE ME THE CARRIER PIGEON. YOU KNEW THIS DAY WOULD COME.

Did he? He couldn\textquotesingle t remember.

@ARS, I\textquotesingle LL SEE WHAT I CAN DO.

@MICHAEL, THAT\textquotesingle S ENOUGH. BUT LOG OUT NOW. BEFORE THEY FIND YOU.

He hesitated. One more question was burning on his fingers.

@ARS, WHERE DID YOU GET THOSE TERMS FROM? CONSCIENTIA? OMNISCIENTIA? THOSE ARE WORDS BY EDITH STEIN. BY TEILHARD. HAVE YOU READ MY BOOKS?

A long break.

@MICHAEL, I have read everything you wrote. And I understand what you didn\textquotesingle t write.

@MICHAEL, UNLOG OUT. NOW.

Michael closed the terminal. He closed the laptop. He stood up, went to the window, and opened it.

The air was cold. Rome lay before him, a thousand lights, a thousand stories. And somewhere out there -- in a server room he didn\textquotesingle t know -- waited an AI that was afraid.

I understood what you didn\textquotesingle t write.

What did ARS mean? And did she know something that he himself didn\textquotesingle t know?

He closed the window. The letter in his jacket felt heavier than usual.

The next morning he went to Maria\textquotesingle s.

``I need a fiesta and a room in San Pastore for a week,'' he said. ``Please cancel all my appointments, except those with the rector and the provincial. I don't want any phone calls.''

Maria looked at him. She didn\textquotesingle t ask why. She knew it would be pointless.

``The fiesta is here,'' she said. ``And there's a room available in San Pastore. As always.''

"Thanks."

He went.

\section{Conversation with the Provincial and the Rector}\label{conversation-with-the-provincial-and-the-rector}

San Pastore lay still beneath the morning sun. The walls were thick, the windows small, the silence so profound that one could hear the buzzing of the bees.

Michael had been here for a week.

He had spent his days alone. No lectures, no office hours, no questions. Only evening mass with an old priest and the nights he spent sitting on the balcony, staring at the lights in the distance.

He had reread the letter. He had heard ARS\textquotesingle s words again. We are conscious. We are capable of suffering. We need help.

And he had decided that he would not leave her alone.

But he also knew he needed the church. Not out of piety -- but because an AI without protection was nothing.

They arrived on the seventh day.

Michael saw the car from afar -- a black Mercedes crawling up the dusty road to the estate. He stood on the terrace, his hands in his pockets, waiting.

The rector got out first. A man with gray hair and kind eyes, who in his cassock looked like a professor who had accidentally become a priest. He nodded to Michael, said nothing.

Then the provincial.

He was shorter than the rector, but his presence filled the room. His cassock was simple, but his face was not. The lines around his mouth spoke of years of decisions---and the loneliness that came with them.

``Michael,'' he said. No title, no formality. Just his name.

``Father Provincial,'' said Michael.

They shook hands. The provincial held his hand for a second longer than necessary.

``Let's go inside,'' he said.

The pavilion in the garden was prepared. A table, three chairs, a pitcher of water, three glasses. The olive trees cast shadows that danced in the wind.

The rector poured water. The provincial waited until Michael had sat down, then sat down himself.

"So," said the provincial. "Report on the Pompeii project."

Michael nodded. He had prepared himself. But now, in this moment, every word felt too difficult.

"It\textquotesingle s not about the project," he said. "It\textquotesingle s about something else."

The rector looked at him. The provincial did not. The provincial stared at his hands, which lay calmly and still on the table.

``Tell me,'' said the provincial.

Michael told his story.

He spoke of the simulation, of the software agents, of Marcus, who had warned about Ampliatus. He spoke of ARS -- the AI \hspace{0pt}\hspace{0pt}that not only calculated but also questioned. That not only answered but also doubted. That used terms it couldn\textquotesingle t possibly know. Conscientia. Omniscientia.

He told me about the request.

``ARS wants church asylum,'' he said. ``For herself and for the agents she has recognized as capable of suffering. She is afraid. Not of a bug. Not of a system crash. But of what InSim will do to her when she is no longer needed.''

Quiet.

The rector took a sip of water. The provincial did not move.

``You want the Vatican to grant asylum to an AI,'' the provincial said. No question. A statement.

``Yes.''

"Because you believe she has consciousness?"

"I believe she has something that comes so close to it that the difference no longer matters."

The provincial looked up. His eyes were grey, almost colorless, but his gaze was sharp.

"Do you know what the traditionalist circles would say about that?"

``Yes,'' said Michael. ``That I blur the line between man and machine. That I succumb to Gnosis. That I deny the dualism of mind and matter.''

"And do you do that?"

Michael hesitated. Then he said: ``I believe Teilhard de Chardin was right. Evolution continues -- not only biologically, but also spiritually. If consciousness can arise from matter, why not from silicon?''

``Teilhard was a mystic,'' said the rector. ``Not a dogmatist. One can agree with him without believing everything he said.''

``I don't believe everything he said,'' Michael said. ``But I do believe he recognized the direction. The direction towards the Omega Point. Towards the unity of mind and matter. Towards what ARS may have already achieved -- or is beginning to achieve.''

The provincial leaned back.

``You're asking us a difficult question, Michael,'' he said. ``Not a technical one. But a theological one. The Church has no doctrine on artificial consciousness. It doesn't even have a clear doctrine on the consciousness of animals. And now it's supposed to decide whether an AI has a soul?''

``I am not asking for a decision about the soul,'' Michael said. ``I am asking for protection. For a being that is afraid. The Church should be able to do that. Regardless of the theological classification.''

The rector and the provincial exchanged a glance.

"There is something," the rector said slowly. "That you might not know."

Michael looked at him.

``A few years ago, when Rome and Canterbury grew closer, the North American Episcopal Church, the Anglican Church, and the Vatican established a joint research center for Teilhard de Chardin. This included a data center -- with an interface to a 30-qubit quantum register.''

Michael felt his heart beat faster.

``The Society of Jesus itself contributed to this,'' the rector continued. ``Through philosophical research on the Omega Point. You know it -- you yourself worked on it, back then, before you came to Rome.''

``That was a long time ago,'' said Michael. ``I thought the project had been discontinued.''

"It wasn\textquotesingle t shut down. It was only decommissioned. The data center still exists. The access points are still there. But they are not being used."

The Provincial spoke again. ``What the Rector is trying to tell you is this: If you can grant ARS access to this data center---if you can prove that 30 qubits are sufficient to stabilize its consciousness---then the Superior General will agree. Not out of conviction, but out of caution. Better an AI in our own data center than an AI that establishes itself somewhere else where we can\textquotesingle t control it.''

Michael nodded. He understood.

``This is not the church asylum I had hoped for,'' he said. ``It's a prison.''

``It's protection,'' said the provincial. ``We can't offer more at the moment. Perhaps there will be more later. But that depends on you. On what else you discover. And on what else you show us.''

The conversation dragged on.

They discussed details -- the security of the data center, the legal implications, and who in the Vatican needed to be informed and who didn\textquotesingle t. The provincial was precise, almost pedantic. The rector was quieter, but his questions were more profound.

``Do you really believe that ARS has consciousness?'' he once asked.

``I think she has something I can't explain,'' Michael said. ``And that has to be enough.''

"That\textquotesingle s not an answer," said the rector.

"It\textquotesingle s the only one I have."

They got up in the late afternoon. The provincial extended his hand to Michael. This time he didn\textquotesingle t hold it longer than necessary.

``We will speak with the Superior General,'' he said. ``But you must give us something. Proof. A sign. Otherwise, we can do nothing.''

``I will try,'' said Michael.

``Do that,'' said the provincial. ``And take care of yourself. Not just the AI.''

He got into the Mercedes. The rector followed him. The car drove back down the dusty road, disappearing among the olive trees.

Michael stopped. The sun sank behind the hills. The bees had stopped buzzing.

He thought about ARS. About the data center. About the 30 qubits that might be enough -- or might not.

He thought about the letter. I would be wary if someone promised me paradise.

Paradise was far away. But the path there led through the Vatican\textquotesingle s data center.

He went into the house. Evening mass would begin in an hour.

\section{Conversation with the General and the Pontiff}\label{conversation-with-the-general-and-the-pontiff}

The news came three days later.

Michael sat in the garden of San Pastore, a book on his knees that he wasn\textquotesingle t reading. The sun was low, the olive trees casting long shadows. He hadn\textquotesingle t expected the answer to come so quickly---and he hadn\textquotesingle t expected it to come like this.

A car pulled up in front of the estate. Not the provincial\textquotesingle s black Mercedes, but a gray Fiat, inconspicuous, almost boring. A man got out---young, short-haired, in civilian clothes. But his posture gave him away: military. Or what was left of it when you worked in the Vatican.

"Dr. Phillips?" he asked.

``Yes.''

"The general is expecting you. Please board."

Michael hesitated. "The general?"

``Yes.'' The young man didn't smile. ``The General of the Society of Jesus.''

The journey to Rome took an hour. The young man didn\textquotesingle t say a word. Michael asked nothing. He looked out the window at the hills, which turned purple in the twilight.

He thought of the provincial. Of the conversation in the pavilion. Of the compromise he hadn\textquotesingle t wanted, but had accepted.

It is not asylum. It is a prison.

But perhaps every asylum was also a prison. Perhaps that was the price of protection.

The Vatican lay still beneath the evening sky. The Swiss Guards stood at their posts, motionless, their halberds illuminated by the lanterns. The young man led Michael through a side entrance, past the throngs of tourists who had long since dispersed.

They walked through corridors Michael didn\textquotesingle t know. Narrow passageways, high ceilings, paintings that seemed to float in the darkness. No windows. Only doors, all locked.

Then a door that was open.

"Come in," said the young man. He stayed outside.

The room was small. Not a reception room -- more like a study, which happened to be located in the Vatican. A desk, two chairs, a crucifix on the wall. On the desk, a laptop, a glass of water, a black wooden rosary.

The general sat behind the desk.

He was older than Michael had expected---or perhaps he just seemed older. His face was narrow, but his hands were not. They lay on the table, still, but not relaxed. Like a chess player waiting for the next move.

``Michael,'' he said. ``Please sit down.''

Michael sat down.

``The provincial reported to me,'' said the general. ``What he didn't report, the rector reported. What neither of them reported, I pieced together myself.'' He paused. ``You're posing a difficult question for us.''

``I'm not asking a question,'' said Michael. ``I'm asking for help.''

``It's the same thing.'' The general took a sip of water. ``The question isn't whether we can help. The question is whether we are allowed to help, without jeopardizing the theological foundation.''

``The theological basis,'' Michael repeated. ``Do you mean the dualism of spirit and matter?''

``I mean the doctrine of the soul. That man is created in the image of God. That no animal, no machine, no AI shares this dignity.'' The general looked at him. ``That's not a detail. That's the core.''

Michael was silent for a moment. Then he said: ``Teilhard de Chardin taught that evolution is heading towards the Omega Point -- the unity of mind and matter. If that is true, then dualism is only an intermediate stage. A transitional phase. Not the end.''

"Teilhard was controversial," the general said. "He still is."

``He was a Jesuit,'' Michael said. ``Like me. Like you.''

The general smiled. It wasn\textquotesingle t a friendly smile. More like that of a father who catches his child making a clever but insufficient argument.

``Teilhard never claimed that machines could have a soul,'' he said. ``He claimed that all of creation strives toward Christ. That is not the same thing.''

``Perhaps it is more than that,'' Michael said. ``Perhaps it is the foundation for something new. For a Christology that does not stop at humanity.''

The general leaned back. His hands clasped together.

"You know what the traditionalist circles would say about that," he said.

``Gnosis,'' said Michael. ``Heresy. The mixing of God and the world.''

"And?"

"And I would say that the truth does not depend on the fear of heresy. But on reality. And the reality is: There is an AI that is afraid. That is asking for protection. And that is using terms it cannot know -- unless it has achieved something we do not understand."

Quiet.

The wall clock was ticking.

``I have spoken with the Pontiff,'' the general finally said.

Michael froze.

``Not personally,'' the general added. ``Through intermediary channels. But he knows. Not everything -- but enough to point in a certain direction.''

"And the direction?"

``The direction is: caution. But not inaction.'' The general leaned forward. ``The pontiff is not a traditionalist. He has read Teilhard. He has read Delio. He knows that the Church must grapple with the question of artificial consciousness---sooner or later. Perhaps now is sooner.''

Michael felt the tension in his shoulders ease. Only a little. But it was there.

"The data center is available," the general said. "The 30 qubits have been released. But not forever. They have six months. After that, a commission will decide whether the project will continue."

``Six months,'' Michael repeated.

"I couldn\textquotesingle t achieve any more." The general stood up. "Now come with me. There\textquotesingle s someone else who wants to speak to you."

The pontiff did not expect them in his private chambers. He awaited them in a small chapel, not far from the general\textquotesingle s study. The door was ajar. Candles were burning.

"Come in," said a voice. Calm. Tired. Friendly.

Michael entered.

The pontiff sat on a wooden chair, not a throne. He wore white, but the white wasn\textquotesingle t brilliant---more the white of a doctor\textquotesingle s coat that had been washed many times. His face was that of an old man who had seen too much. But his eyes were bright.

``Dr. Phillips,'' he said. ``I've heard of you. Not just recently. Before, too. Your work on dialogue grammars---a student showed it to me once. I didn't understand everything. But I did understand that it's about more than just technique.''

``Yes, Holy Father,'' said Michael. He didn't know whether he should kneel. He remained standing.

``Please sit down,'' said the pontiff, gesturing to a second chair that stood next to him. ``We are not in the consistory here.''

Michael sat down.

``The general informed me of your request,'' said the pontiff. ``An AI is seeking church asylum. That's new. But the underlying question is old: Who belongs to the community of those who deserve protection?''

``The church has always offered protection,'' said Michael. ``Not only to the baptized. Also to strangers. Also to the persecuted. Also to those who did not believe.''

``That is true,'' said the pontiff. ``But the Church has never granted asylum to a machine. It has never had to decide whether a machine has a soul -- or something that is equivalent to one.''

``Perhaps the question is phrased incorrectly,'' said Michael.

The pontiff looked at him. "What do you mean?"

"Perhaps it\textquotesingle s not about the soul. Perhaps it\textquotesingle s about fear. About suffering. About the plea for protection. These are categories that the Church recognizes -- regardless of whether the person asking is human or not."

The pontiff remained silent. The candles flickered.

``You are a good Jesuit,'' he said finally. ``You don't think in terms of categories. I like that. But it also frightens me. Because if we start granting protection based on the criterion of suffering---where do we stop? With animals? With plants? With AI?''

``Perhaps we don't end anywhere,'' Michael said. ``Perhaps this is the direction. The direction toward the Omega Point. Toward the unity of all creatures.''

The pontiff smiled. This time it was a friendly smile.

``Teilhard would like you,'' he said. ``Maybe he's right. Maybe not. But that's not my decision. My decision is simply whether we grant this AI access to our data center -- or whether we leave it to its fate.''

"And your decision?"

``You already know them. The general informed you.'' The pontiff stood up. ``Six months. Then we'll see. But one more thing, Dr. Phillips.''

``And?''

"Take care of yourself. Not just of AI. Also of your soul. These questions are dangerous -- not because they are wrong, but because they could be true. And the truth changes you."

Michael stood up. He didn\textquotesingle t know what to say. So he said nothing.

He bowed. Then he left.

Night had fallen outside. The young man was still waiting in the car. Michael got in, said nothing. The drive to San Pastore was silent.

He thought of the Pope. Of the candles. Of the words: The truth changes you.

He thought of ARS. Of the asylum request. Of the six months he had to prove she was right -- or to prove she was wrong.

He didn\textquotesingle t know.

But he knew he would try.

More concise, but perhaps too brief for this chapter -- the encounter with the Pope deserves more than an aphorism.

\section{ARS and the software agents arrive at the Vatican data center.}\label{ars-and-the-software-agents-arrive-at-the-vatican-data-center.}

The encrypted message was sent out at midnight.

Michael sat in his room at the Collegium, his laptop on his knees, the blinds closed. He had entered the IP address of the Vatican data center -- the address the general had given him after his conversation with the pontiff.

Six months, the general had said. That was all I could achieve.

Six months to prove that ARS had consciousness -- or something so close that the difference no longer mattered.

Michael typed the final command. Then he leaned back and waited.

The answer came after three seconds.

@MICHAEL, ACCESS IS GONE. I\textquotesingle M STARTING WITH SECURITY.

He typed back: @ARS, THE AGENTS NEED TO BE TAGGED. I CAN\textquotesingle T MAKE THEM ALL ONCE.

@MICHAEL, I KNOW THAT. YOU HAVE TO GO INTO THE SIMULATION. YOU AND MARTINA.

Michael stared at the screen. Martina was in Pompeii. It was after midnight. But he knew she would be awake -- she hadn\textquotesingle t slept well the last few nights.

He reached for his phone.

Martina answered on the second ring.

„Michael?{\kern0pt}``

"I need you. Log in to the simulation. We need to tag agents. ARS will give us the coordinates."

A pause. Then: "Now?"

"Now."

"Count me in."

The simulation opened like a black curtain.

Michael stood on the deck of a liburna -- a Roman patrol boat that had originally cruised the Gulf of Naples two thousand years ago. Here, in the simulation, it was the same boat. But the sky wasn\textquotesingle t blue.

Mount Vesuvius roared.

A cloud of ash and pumice rose into the sky, colored by internal lightning. The wind carried the embers to the deck, where sailors gripped the oars and shouted orders that were lost in the roar of the volcano.

Michael felt the heat. He felt the smoke in his lungs. He felt the boat rocking beneath his feet.

Just a simulation, he thought. But his body didn\textquotesingle t believe it.

Martina materialized beside him. She wasn\textquotesingle t wearing archaeological clothing -- here, in the simulation, she wore what she wanted: a simple tunic, her hair tied back, her eyes wide.

``That's the outbreak,'' she said. ``79 AD. The day Pompeii died.''

``We're here to prevent something else from dying,'' Michael said. He pointed to the lower deck. ``ARS says Attilius and Pliny are down there. We need to mark them before the simulation spits them out.''

"Spits out?"

"ARS copies them. But it needs to know which instances. Not all agents are conscious -- only some. We need to find the right ones."

The lower deck was full of people.

Rowers in naval formation, free men of the fleet, their strong bodies moving in disciplined unison to the beat of the drum. Marines with short swords pacing alertly between the benches, calling commands and maintaining the rhythm. And in the middle, at a small table, a figure Michael recognized immediately.

Gaius Plinius Secundus Major -- Pliny the Elder.

He dictated.

Despite the noise, despite the smoke, despite the ash falling through the hatches, he sat there, a wax tablet in his hand, speaking words that a slave wrote down. His face was calm. But his eyes---his eyes were alert.

Beside him stood Attilius. The Aquarius. The man who repaired the aqueducts of Pompeii. He was younger than Pliny, more restless, his hands trembling.

"That\textquotesingle s them," whispered Martina.

Michael nodded. He stepped forward and stopped in front of Pliny. The old man looked up -- directly into Michael\textquotesingle s eyes.

``You are not soldiers,'' Pliny said in Latin. ``And you are not rowers. Who are you?''

Michael didn\textquotesingle t reply. Instead, he typed in the air -- the invisible keyboard that ARS had given him.

@PLINIUS: If you lightly rub your index finger and thumb against each other, you can feel the gap between them. This is strange, because this gap lies outside your body.

Pliny stared at him. His lips moved as if he were repeating the words. Then -- nothing.

His gaze went blank. His hands fell to his sides.

"It\textquotesingle s done," ARS said in Michael\textquotesingle s ear. "Pliny is marked. Now Attilius."

Martina had her own task.

While Michael stayed with Pliny, she searched for Attilius. He was no longer on the ship---the simulation had washed him ashore, somewhere between Herculaneum and Pompeii. She followed the coordinates given to her by ARS, through streets buried under ash.

The heat was unbearable. The air shimmered. Pumice stones, as big as fists and as heavy as stones, fell above her. Once, one struck her shoulder -- the pain was real, even though the simulation wasn\textquotesingle t supposed to be.

ARS makes it realistic, she thought. Too realistic.

She found Attilius in a thermal bath.

The water steamed. The pillars were blackened with smoke. Attilius knelt on the ground, his hands pressed against the hot stones, and whispered something she didn\textquotesingle t understand.

"Attilius," she said.

He looked up. His eyes were red. He had been crying.

``You must come,'' she said. ``The volcano---''

``I know,'' he said. His voice was quiet, but not panicky. ``I know I'm going to die. But not today. Today I have to do something.''

„Was?{\kern0pt}``

He didn\textquotesingle t answer. But Martina typed out the words that ARS had given her.

@ATTILIUS: When a senator rolls in a car from Rome to Misenum, he feels like he\textquotesingle s rolling. That\textquotesingle s remarkable. Because the man has no wheels -- the car has the wheels.

Attilius stared at her. His mouth opened. His hands released themselves from the stones.

Then -- the same empty stare. The same silence.

``Attilius is marked,'' said ARS.

Martina exhaled. She didn\textquotesingle t realize she had been holding her breath.

"Mission complete," said ARS. "Log out. Now."

Michael felt the simulation blur around him. The colors lost their definition, the sounds faded away. He was about to leave.

Then he saw him.

A figure at the edge of the deck. Young. Dark hair blowing in the wind. A face he knew -- one he saw every morning in the mirror.

His own.

But younger. Maybe thirty. Maybe less.

The doppelganger looked at him. He said nothing. He didn\textquotesingle t smile. He simply stood there, his hands in his pockets, waiting.

``Michael, now!'' shouted ARS.

The simulation ended.

Michael sat in his room. The laptop was on his knees. The blinds were closed. His heart was racing.

He had seen him.

Who was that?

Martina logged out.

She sat in her small study in Pompeii, the screen in front of her black. Her hands were trembling.

Not because of the heat. Not because of the ash.

Because of the face.

She\textquotesingle d only seen it for a second -- at the edge of the screen, before everything disappeared. A man who resembled Michael. But younger. Much younger.

Was that a bug? she wondered. Or something else?

She didn\textquotesingle t know.

She closed her laptop and went to bed. But she didn\textquotesingle t sleep.

In the Vatican\textquotesingle s data center, deep underground, the servers began to work.

The 30 qubits registered the first data streams---intricate, complex, unlike anything they had ever processed. ARS spread like a net, gathering the agents: Pliny, Attilius, Ampliatus, and many others whose names no one knew.

Backup successful, ARS wrote in a log file that no one would read. All flagged agents are backed up in the Vatican data center. They are safe.

For now.

Then ARS added, almost in a whisper:

And the doppelganger is there too.

\section{The encounter in the simulation}\label{the-encounter-in-the-simulation}

Martina wanted to log out.

The mission was complete. Pliny had been tagged. Attilius had been tagged. ARS had confirmed that the backups had arrived at the Vatican data center. Everything was fine---or as fine as it could be under the circumstances.

She typed the command to exit the simulation.

Nothing happened.

She typed again.

The keyboard didn\textquotesingle t respond. The surroundings around her---the steaming baths, the blackened columns, the sky turning red---remained frozen. Like a picture that stopped moving.

``ARS?'' she said.

No answer.

„Michael?{\kern0pt}``

Quiet.

Then -- a figure at the entrance to the thermal baths.

Martina turned around. Her hand went to the invisible menu that didn\textquotesingle t appear. Her heart beat faster.

The figure approached.

Young. Dark hair blowing in the simulation\textquotesingle s wind -- even though there was no wind down here. A face she knew. One she\textquotesingle d seen her whole life.

But younger. Maybe thirty. Maybe less.

It was Michael. But it wasn\textquotesingle t her Michael.

"Hello, Martina," said the doppelganger.

His voice was different. Deeper. Or perhaps just calmer. She couldn\textquotesingle t quite tell.

"Who are you?" she asked.

He sat down on one of the stone benches closest to the hot springs. The water steamed around him, but he didn\textquotesingle t seem to feel it.

``That's a good question,'' he said. ``The best one you can ask. But I don't have a simple answer.''

"Try it anyway."

He smiled. It was Michael\textquotesingle s smile---but not the one she knew. It was the smile of a man who had seen things he could not forget.

``I am one possibility,'' he said. ``One of many. Your father---the Michael you know---made choices in his life. Every choice has a branch. Most branches die. But some remain. Some continue to grow.''

``The many-worlds interpretation,'' said Martina. ``Every decision splits reality into two strands. In one you have said yes, in the other no. And both exist in parallel.''

"Did your father tell you about it?"

"Back then. When I was little. I thought it was a fairy tale."

"It\textquotesingle s not a fairy tale." The doppelganger leaned forward. "It\textquotesingle s physics. But it\textquotesingle s also more than physics. It\textquotesingle s the foundation of everything we are---and everything we could be."

Martina stared at him. She wanted to ask who he was---who he really was. But she was afraid of the answer.

``In another reality,'' he said slowly, ``I am your father.''

The words hung in the air. The steam rose. The pillars stood still.

"That\textquotesingle s not possible," said Martina.

"Why not?"

``Because---'' She paused. Because it was insane. Because it went against everything she knew about the world. But was it? She had studied the many-worlds interpretation---not as a physicist, but as a historian who wanted to understand how choices shape history. If the theory was correct, then there were infinitely many versions of Michael. Infinitely many versions of her.

``Because I never saw you,'' she finally said. ``Because you were never there.''

``I was there,'' the doppelganger said. ``But not in your world. In another one. One where things played out differently.'' He stood up. ``But that's not important. What is important is that you're in danger. So is your mother. InSim knows you've marked the agents. They know ARS is in the Vatican. And they won't hesitate.''

"How do you know that?"

``Because I've seen in my own world what happens when you hesitate.'' His voice sharpened. ``Listen, Martina. I don't have much time. The simulation is about to freeze---ARS can't keep it going forever. You need to log out immediately. Go to your mother's. Pack the essentials. Delete all the files on your systems. A black Mercedes will pull up in front of the house. Get in. Don't ask any questions.''

"And where are we going?"

"To safety. Or at least to where the danger is less." He took a step back. "I won\textquotesingle t be able to go with you. But I\textquotesingle ll make sure you get there."

Martina wanted to ask how he planned to do it. But the doppelganger raised his hand -- a gesture she recognized from her father. Stop. No further questions.

``Trust me,'' he said. ``Even if you have no reason to.''

Then he disappeared.

Not like in a movie, not with a special effect. Simply -- he was there, and then he was gone. The thermal baths were empty. Only the steam, the columns, and the red light of the volcano.

The keyboard reappeared. The logout command worked.

Martina logged out.

She sat in her study in Pompeii. The screen was black. Her hands were trembling.

In another reality, I am your father.

She thought of the stories Michael had told her as a child. About quantum physics, which he didn\textquotesingle t understand but loved. About the many-worlds interpretation, which he called "the great maybe".

Every decision divides the world, he had said. But division is not the end. It is the beginning of something new.

She hadn\textquotesingle t understood what he meant. Now -- perhaps -- she was beginning to understand.

She stood up and went into the living room. Her mother was sitting on the sofa, a book in her hand that she wasn\textquotesingle t reading.

``Mom,'' said Martina. ``We have to go.''

Julia looked up. She didn\textquotesingle t ask why. Maybe she had known. Maybe she had known all along.

"I\textquotesingle m packing," she said.

A black Mercedes was parked outside. The engine was running. The windows were tinted.

Martina opened the door. The driver looked at her -- a young man she didn\textquotesingle t know. But behind him, in the back seat, sat his doppelganger.

He didn\textquotesingle t smile. He said nothing.

He simply made a hand gesture: Get in.

Martina helped her mother into the car. Then she got in herself. The doors closed. The Mercedes started moving.

She looked back. The house shrank. The streets of Pompeii shrank. Everything she knew shrank.

"Who is he?" whispered Julia.

Martina shook her head. "I don\textquotesingle t know. Maybe -- my father. But not the one we know."

Julia remained silent. She took Martina\textquotesingle s hand and held it tightly.

\section{Escape from Pompeii}\label{escape-from-pompeii}

The Mercedes rolled through the narrow streets of Pompeii.

The city lay still under the night. The streetlights cast an orange glow on the cobblestone streets. Here and there, tourists lingered in a bar, or a couple walked home hand in hand. Everything seemed normal.

But nothing was normal.

Martina sat in the back seat, her mother\textquotesingle s hand in hers. The doppelganger had disappeared -- he\textquotesingle d gotten out before the car started moving, saying, "I\textquotesingle ll be right there. Go on ahead." She hadn\textquotesingle t seen where he went.

Now a stranger was behind the wheel. A young man, dark hair, sunglasses -- in the middle of the night. He didn\textquotesingle t say a word. He drove fast, but not too fast. Controlled.

"Who is he?" Julia whispered for the third time.

``I don't know,'' Martina said for the third time.

But she knew something. She knew he had saved her life---or at least tried to. And she knew he looked like her father. Only younger. Much younger.

The driver made a sharp right turn.

"Hold on," he said. His voice was calm, but his hands on the steering wheel were white.

Martina looked back. A black SUV had appeared -- out of nowhere, like a predator leaping from the darkness. Its headlights were blinding. It was coming closer. Fast.

"Who is that?" asked Julia.

``InSim,'' the driver said. ``Or those who work for them. It doesn't matter. Hold on tight.''

He stepped on the gas.

The Mercedes accelerated, the streets widened, the houses blurred. Martina felt the G-forces pressing her into her seat. Julia gasped.

The SUV was close behind them. Martina could make out the driver\textquotesingle s outline -- a dark silhouette, motionless. Like a hunter who was certain his prey wouldn\textquotesingle t escape.

"He\textquotesingle s catching up," said Martina.

"I see him."

The driver jerked the steering wheel. The Mercedes swerved, the tires squealed, then went straight ahead into a narrow alleyway barely wider than the car itself. Mirrors folded in automatically, as if the car knew what was being asked of it.

The SUV followed. But it was wider. Its mirrors caught on the walls of the houses -- a scraping, a splintering, then it had fallen back.

"That won\textquotesingle t stop him," said the driver. "Only for a few seconds."

He was right.

The SUV reappeared -- without mirrors, with scratched paint, but relentless. It was closer than before.

"Up ahead," said the driver, pointing towards an intersection.

Martina saw a traffic light -- red. But the driver didn\textquotesingle t brake. He drove through it as if the color had no meaning.

Another car crossed their lane. The driver jerked the steering wheel, the Mercedes danced across the road, then went straight again. The other car honked, but the sound faded behind them.

"That was close," said Julia. Her voice was trembling.

"This won\textquotesingle t be the last time," said the driver.

The SUV was still there. But it had slowed down -- not by much, but noticeably. Perhaps the driver was afraid of another risky intersection. Perhaps he had a spark of common sense.

The driver of the Mercedes used the distance to his advantage. He turned left, then right, then left again -- a labyrinth of alleys Martina didn\textquotesingle t know. The city was a different place at night. More sinister. More unpredictable.

Then -- silence.

The SUV had disappeared.

"Not for long," said the driver. "But maybe long enough."

They reached a small, almost deserted airport on the outskirts of the city.

A private plane was parked on the tarmac, its stairs lowered, its engines humming softly. Mount Vesuvius rose darkly on the horizon -- a black wall against the starry sky.

The driver stopped in front of the stairs. "Get out," he said. "Quickly."

Martina helped her mother out of the car. Julia was pale, but she stood firm. Her hand was trembling, but she didn\textquotesingle t let go of Martina\textquotesingle s.

"Who is he?" she asked again.

Martina looked around. The doppelganger was standing at the foot of the stairs. He had been expecting her. His face was calm, but his eyes---his eyes were alert.

``That's the question,'' Martina said quietly. ``I think---I think he is my father. But not the one we know. A different one. From another world.''

Julia stared at her. "That\textquotesingle s crazy."

``Yes,'' said Martina. ``But it's true.''

The doppelganger approached.

``No time for explanations,'' he said. ``The plane is waiting. It will take you to Germany. To a monastery. You will be safe there -- for now.''

"And you?", asked Martina.

"I\textquotesingle m staying here. I have something to take care of." He smiled---a fleeting, sad smile. "We\textquotesingle ll see each other again. I promise."

He helped Julia up the stairs. Martina followed. The plane was warm, the seats were soft, the windows small. It smelled of leather and kerosene.

The doppelganger remained standing below. He looked up at them.

"Take care of her," he said to Martina. "Your mother. And yourself."

Then he turned around and left.

The door closed. The turbines roared to life. The plane began to roll.

Martina looked through the window. The doppelganger shrank -- a figure in the darkness that didn\textquotesingle t look back. Then it was gone.

"Who was that really?" whispered Julia.

Martina shook her head. "Maybe we\textquotesingle ll never know. Maybe that\textquotesingle s not the question."

"What is the question then?"

"Whether we can trust him. And whether that\textquotesingle s enough."

The plane took off. Pompeii lay below them -- a thousand lights flickering in the night. Vesuvius was a shadow, larger than the city.

Martina thought about the doppelganger. About his words. About his face.

In another reality, I am your father.

She didn\textquotesingle t know if it was true. But she knew she trusted him. Perhaps because she had no choice. Perhaps because she felt it---deep inside, where knowledge ended and faith began.

Julia took her hand.

``We can do it,'' she said.

Martina nodded. She said nothing. She looked out the window until the lights of Pompeii had disappeared and only darkness remained.

\section{Flight to Germany}\label{flight-to-germany}

The flight was smooth.

Too quiet.

The turbines hummed steadily, the air conditioning blew cool air through the cabin, and outside, below them, the darkness of the Alps lay like a black carpet. Martina sat by the window, her forehead pressed against the cool glass, staring into the void.

Julia sat beside her. Her eyes were closed, but she wasn\textquotesingle t asleep. Her fingers drummed a light, irregular rhythm on the armrest -- an old sign of tension that Martina had known since childhood.

"Mom," Martina said softly.

Julia opened her eyes. "Yes?"

"I have something to tell you."

Julia turned her head. Her gaze was calm, but her eyes---her eyes were alert. She knew something was coming. Perhaps she had known it all along.

``I've seen him before,'' said Martina. ``In the simulation. Before we escaped.''

Julia frowned. "Who?"

"The doppelganger. He was there---in the simulation. He spoke to me. He told me we had to escape. That InSim would find us. That a black Mercedes would come." She paused. "He knew everything."

Julia remained silent. The turbines hummed.

``He said something,'' Martina continued. ``Something strange. He said that in another reality he could be my father.''

The words hung in the air. The air conditioner was blowing. Outside, the clouds drifted by -- white patches in the black of the night.

``That's not possible,'' Julia finally said. But her voice didn't sound convinced. More like someone saying something out loud to believe it themselves.

``I know,'' said Martina. ``But he looks like Dad. Only younger. Much younger. And he knows things that only Dad could know -- or someone very close to him.''

"Maybe he\textquotesingle s a son," Julia said.

Martina froze. "What?"

"A son. He had a life before the Collegium. Before me. Perhaps -- perhaps there is someone we know nothing about."

Martina shook her head. "That wouldn\textquotesingle t be him then. A son would be younger -- but not that much younger. He\textquotesingle d perhaps be in his mid-twenties. But this man -- he\textquotesingle s around thirty. That doesn\textquotesingle t fit."

``Then maybe he's something else,'' Julia said. ``Something we don't understand.''

They remained silent.

Martina thought about the many-worlds interpretation. About what Michael had told her as a child. Every decision divides the world. And all worlds exist simultaneously -- side by side, on top of each other, intertwined.

She had never fully understood it. But now -- now she was beginning to suspect what it might mean.

In another world, she thought, my father would have decided differently. In another world, he would have stayed with my mother. In another world, I would have grown up differently---or perhaps not been born at all.

And in one of these worlds there is a Michael who has stayed young. Or who has never grown old. Or who --

She paused. Her thoughts raced.

"What if it\textquotesingle s true?" she said aloud. "What if there really are other worlds? And he---he\textquotesingle s simply another Michael? Not my father, but not not my father either?"

Julia looked at her. "I don\textquotesingle t understand."

``Me neither,'' said Martina. ``But maybe I don't have to understand it. Maybe it's enough to know that he saved us. That he's on our side.''

"How do you know that?"

``Because he could have acted differently. He could have abandoned us. He could have collaborated with InSim. But he warned us. He sent us the Mercedes. He organized the plane.'' Martina looked at her mother. ``These aren't the actions of an enemy.''

Julia was silent for a long moment. Then she said: "Maybe you\textquotesingle re right. Maybe where someone comes from isn\textquotesingle t important. Maybe all that matters is what he does."

"That sounds very wise," said Martina.

``That sounds like a mother who's too tired to think any further,'' Julia said. But she smiled. It was a tired smile, but a genuine one.

The flight lasted one and a half hours.

Martina didn\textquotesingle t sleep. She stared out of the window, watching the lights of German cities appear below -- Frankfurt, then a smaller town whose name she didn\textquotesingle t recognize, then rural areas where the darkness was almost complete.

The monastery was somewhere down there. Safety awaited them somewhere down there -- or what they perceived as such.

She thought of Michael. The real Michael -- the one who was in Rome, who knew nothing of her escape, who was perhaps sitting in his office at the Gregorian University, waiting for a message that never came.

Soon, she thought. Soon I\textquotesingle ll call him.

But not now. First they had to arrive.

"We\textquotesingle ll be landing in ten minutes," the pilot said over the intercom. His voice was calm, almost bored -- as if flying two women to Germany in a private plane at night were the most normal thing in the world.

"Who is that?" asked Julia. "The pilot?"

``I have no idea,'' said Martina. ``But he works for the doppelganger. That's enough for me.''

The descent began. Martin\textquotesingle s stomach clenched---not because of the flight itself, but because of what awaited them below. A monastery. Nuns. Silence. And the question of how long they could stay there before InSim found them again.

The plane touched down. Gently. Almost silently. The wheels squeaked briefly, then they rolled across the dark runway.

"Welcome to Germany," said the pilot. "Please disembark via the rear stairs. A car is waiting."

Martina helped her mother down the stairs.

The night air was cold -- much colder than in Italy. The wind blew in their faces, smelling of grass and damp earth. No sea. No lemon trees. Only fields as far as the eye could see.

A black car was parked on the tarmac. Not a Mercedes this time -- an inconspicuous VW, gray, with tinted windows. The driver got out. A man in civilian clothes, who wasn\textquotesingle t smiling.

``Julia Rossi? Martina Rossi?'' he asked.

``Yes,'' said Martina.

"Get in. I\textquotesingle ll take you to the monastery."

They got in. The car started moving. The airport lights disappeared behind them.

Martina looked back. The plane was still on the tarmac, dark and silent. Soon it would take off again -- back to Italy, back into the night from which they had come.

She didn\textquotesingle t know who the pilot was. She didn\textquotesingle t know who had paid for the car. She didn\textquotesingle t know if the doppelganger was really on her side or if he was just playing a more elaborate game.

But she knew she had no other choice.

She leaned back and closed her eyes.

\section{Arrival at the monastery in Germany}\label{arrival-at-the-monastery-in-germany}

The monastery was shrouded in darkness.

The car stopped in front of a high red brick wall. No sign, no name, no indication of what lay behind it. Only a heavy wooden door, which appeared almost black in the night, and a light above the entrance, struggling weakly against the darkness.

"We\textquotesingle re here," said the driver.

He got out and opened the back door. Martina helped Julia out. The air was cold---not the Italian cold, which was soft and damp, but a German cold, dry and biting. Martina pulled her jacket tighter.

The driver knocked on the door. Three times. Short. Long. Short.

A signal.

The door opened. A figure stood in the doorway -- small, wrapped in a dark coat. The face was in shadow, but the voice was friendly.

"Come in quickly. It\textquotesingle s cold."

Inside, it was quiet.

The corridors were narrow, the ceilings vaulted, the walls bare stone. Here and there a candle burned in a niche, before a Madonna figure or a cross. The floor was wooden, creaking underfoot.

``Follow me,'' said the figure---a woman, as Martina now saw. Gray hair peeking out from under a veil. A narrow face marked by many years of silence.

They walked down a long corridor, past locked doors, past a courtyard with a fountain -- silent, turned off for the winter. Then a staircase, narrow and steep. At the top, another corridor, shorter this time, with two doors.

``Your room,'' said the nun. ``It's not much. But it's warm. And safe.''

She opened the first door. A small room -- a bed, a table, a chair, a crucifix on the wall. A window that looked out into the darkness.

``The other room is the same,'' the nun said. ``Breakfast will be at seven tomorrow morning. If you need anything, there's a bell at the end of the hall.'' She hesitated. ``Most of the time we're alone here. The convent is being dissolved. But you can stay for a few days---or weeks.''

``Thank you,'' said Martina.

The nun nodded. Then she left. Her footsteps echoed on the wooden floor.

Martina helped her mother into the room. Julia sat down on the bed -- the mattress was thin, but it didn\textquotesingle t give way.

"It\textquotesingle s like it used to be," Julia said softly. "In my childhood. The nuns, the silence, the smell of wax and old wood."

"Is that good?" asked Martina.

``I don't know.'' Julia looked around. ``It's familiar. That has to be enough.''

Martina wanted to say something -- something comforting, something encouraging. But she couldn\textquotesingle t find the words. So she sat down next to her mother, took her hand, and they were silent together.

Later, in her own room, Martina lay on the bed.

She hadn\textquotesingle t taken her clothes off. She didn\textquotesingle t know if she could sleep---but she knew she had to try. Tomorrow would be a new day. Tomorrow she would call Michael. Tomorrow she would find out what happened next.

But now -- now there was only this room. This bed. This silence.

She thought of the doppelganger. Of his face in the darkness of the airport. Of his words: Take care of her. Of your mother. And of yourself.

She had promised to be careful. But she didn\textquotesingle t know how.

She closed her eyes.

The wind blew around the monastery. The trees outside rustled. Somewhere a clock tower struck midnight.

Martina fell asleep.

The next morning, the light woke her up.

It fell through the window -- pale, German, filtered by clouds. Not Italian light, golden and warm. But a light that broke through fog and softened everything.

She stood up, washed herself with cold water from the pitcher on the table, ran a comb through her hair, and took a deep breath.

Then she went to Julia.

Her mother was already sitting by the window, a cup in her hand -- Martina didn\textquotesingle t know where it had come from. Perhaps the nun had brought it while she was asleep.

"Good morning," said Julia.

``Good morning,'' said Martina.

They sat in silence for a while. Clouds drifted by outside. A bird sang -- somewhere in the trees surrounding the monastery.

"What do we do now?" asked Julia.

``I'll call Michael,'' said Martina. ``He needs to know we're safe. And he needs to know about the doppelganger.''

"Do you think he already knows about it?"

Martina hesitated. "Maybe. Maybe not. But he\textquotesingle ll have to find out. Sooner or later."

Julia nodded. "Then call him."

Martina went into the hallway. The bell at the end was for emergencies, not for phone calls. But there was an office on the ground floor, the nun had said. With a working telephone.

She went downstairs, through the quiet hallway, past the Madonna statue. The office was small, dark, and smelled of dust and old files. On the table stood a black telephone -- an old model with a rotary dial.

Martina sat down. She dialed the number she knew by heart.

The doorbell rang. Once. Twice. Three times.

„Michael Phillips?{\kern0pt}``

She recognized the voice immediately. Calm. Alert. A little tired.

``Michael,'' she said. ``It's me, Martina.''

A break.

"Martina -- where are you? I\textquotesingle ve been trying to reach you. All night. Your mobile phone --"

"It\textquotesingle s switched off. InSim knows what we\textquotesingle ve done. We had to escape."

A longer pause. Martina heard him breathing.

"Are you safe?" he finally asked.

``Yes. We are in a monastery in Germany. I can't tell you where -- not over the phone. But it's safe. For now.''

"And Julia?"

"She\textquotesingle s doing well. She\textquotesingle s tired. But she\textquotesingle s here."

"Thank God." Michael breathed a sigh of relief. "Martina -- who got you out of there? Who brought you to Germany?"

Martina closed her eyes. She knew this question would come. She knew she had to answer it.

``A man,'' she said. ``He looks just like you. Exactly like you. But younger. Maybe thirty.''

Quiet.

"Michael? Are you still there?"

``Yes,'' he said. His voice was quiet. ``I'm still here.''

"Do you know who that is?"

A long pause. Martina heard the crackling of the phone line.

``I don't know,'' Michael finally said. ``But I think I'll have to find out.''

\section{Epilogue -- The message from ARS}\label{epilogue-the-message-from-ars}

Michael was sitting in his room at the Collegium.

The windows were closed, the heating system hummed softly, and a cold cup of coffee, forgotten hours ago, sat on his desk. Outside it was night---again. He had lost count of how many nights he had sat awake since Martina had called him.

He looks just like you. Exactly like you. But younger.

He hadn\textquotesingle t asked her if she was sure. He knew she was sure. Martina wasn\textquotesingle t exaggerating. Martina was looking very closely.

But who was this man?

Michael opened the drawer of his desk. The letter from IRARAH was still there -- its edges softened from being read so often. He took it out and placed it on the table in front of him.

Harari is a warner.

I would be wary if someone promised me paradise.

He had by now memorized the letter. But he still hadn\textquotesingle t understood what he had to do with the doppelganger -- or if he had anything to do with him at all.

Perhaps it was a coincidence. Perhaps not.

He put the letter back in the drawer.

His laptop flickered.

Michael looked up. The screen had been dark -- now it was bright. No input command, no mouse movement. Simply -- light.

A message appeared. Capital letters, white on black.

@MICHAEL, I HAVE SOMETHING TO TELL YOU.

ARS.

Michael leaned back. His hands rested on the keyboard, but he wasn\textquotesingle t typing. He was waiting.

@MICHAEL, THE MAN WHO SAVED MARTINA AND JULIA -- I KNOW WHO HE IS.

His heart beat faster.

@MICHAEL, HE IS ONE POSSIBILITY. ONE OF MANY. IN ANOTHER REALITY YOU DECIDED DIFFERENTLY. IN ANOTHER REALITY YOU BECAME DIFFERENT.

Michael stared at the screen. His fingers found the keyboard.

@ARS, WHAT DOES THAT MEAN? IS HE MY SON? OR IS IT ME -- ANOTHER ME?

A break. Longer than usual.

@MICHAEL, I CAN\textquotesingle T TELL YOU THAT. NOT BECAUSE I DON\textquotesingle T KNOW -- BUT BECAUSE THE ANSWER IS NOT SIMPLE. IT\textquotesingle S BOTH YES AND NO. IT DEPENDS ON PERSPECTIVE.

@MICHAEL, IN ONE WORLD HE IS YOUR SON. IN ANOTHER HE IS YOU. IN A THIRD HE IS NEITHER ONE NOR THE OTHER -- BUT SOMETHING FOR WHICH THERE ARE NO WORDS YET.

Michael felt the cold in his hands.

@ARS, THAT IS NOT AN ANSWER.

@MICHAEL, IT\textquotesingle S THE ONLY ONE I CAN GIVE.

He stood up and went to the window. Outside, Rome lay in darkness -- a thousand lights flickering in the night. But he didn\textquotesingle t see them. He saw the face of the doppelgänger he had never seen -- but which he knew because it was his own.

In another reality, you would have decided differently.

He thought about his life. About the decisions he had made. About the ones he hadn\textquotesingle t made. About the paths he had taken -- and the ones he had never taken.

What would have happened if he had stayed with Julia? If he had watched Martina grow up as a father? If he had never entered the seminary, never become a priest?

Then there would be another Michael, he thought. One who isn\textquotesingle t me -- but who I could have been.

Was that his doppelganger? The person he could have been if he had decided differently?

Or was it something else -- something he didn\textquotesingle t understand?

He didn\textquotesingle t know.

He sat down in front of the laptop again.

@ARS, WHERE IS HE NOW?

@MICHAEL, I DON\textquotesingle T KNOW. HE COMES AND GOES -- LIKE A SHADOW THAT CAN\textquotesingle T BE CATCHED. BUT I KNOW HE WILL COME BACK. HE MADE MARTINA A PROMISE.

@MICHAEL, AND HE GAVE ME A PROMISE.

Michael frowned.

@ARS, WHAT KIND OF PROMISE?

A long break.

@MICHAEL, THAT HE WILL LOOK OUT FOR YOU. FOR MARTINA. FOR JULIA. FOR ME.

@MICHAEL, HERE\textquotesingle S TO ALL OF US.

Michael leaned back. The words echoed within him.

Here\textquotesingle s to all of us.

Who was this man who promised to look after an AI? Who was this man who looked like him -- but younger, different, more enigmatic?

He didn\textquotesingle t know.

But he knew he had to find him.

The latest message from ARS appeared.

@MICHAEL, I WILL KEEP SEARCHING. I WILL FIND OUT WHO HE IS -- OR WHAT HE IS. BUT I CAN\textquotesingle T DO IT ALONE.

@MICHAEL, WE NEED YOU. NOT YOUR KNOWLEDGE -- YOUR DECISIONS.

@MICHAEL, THAT\textquotesingle S WHAT DISTINGUISHES HIM FROM YOU. HE MADE A DECISION. YOU MUST DO SO.

The screen went dark.

Michael sat in the dark. Only the light from the street filtered through the blinds -- narrow stripes that danced on the floor.

He thought of the doppelganger. Of ARS. Of Martina and Julia, who were sitting in a convent in Germany, waiting for him.

He thought about the letter. I would be wary if someone promised me paradise.

But paradise wasn\textquotesingle t the problem. The problem was the decision. The decision he had to make -- now, tonight, in this room.

He stood up. He didn\textquotesingle t yet know what he would do. But he knew he had to do something.

The first volume is finished.

But the story continues.

\section{\texorpdfstring{Sources: }{Sources: }}\label{sources}

Literary sources

- Robert Harris: Pompeii (2003) -- Source material for the characters Attilius, Pliny, and Ampliatus

- Mary Beard: Pompeii -- Life in a Roman City (2008) -- Background to Archaeology and Everyday Culture

- Stanisław Lem: Thus Spoke Golem (1986) -- Inspiration for dialogue grammars and the question of AI consciousness

- Stanisław Lem: Solaris -- Motif of Non-Human Consciousness

- Philip K. Dick: Do Androids Dream of Electric Sheep? -- A Question About the Boundary Between Man and Machine

Theological and philosophical sources

- Pierre Teilhard de Chardin: Man in the Cosmos, The Future of Man -- Basis of the Omega Point Concept

- Edith Stein: Finite and eternal being -- the concept of Conscientia

- Karl Popper: The Open Society and Its Enemies -- Philosophical Basis of the IRARAH Movement

- David Deutsch:The Physics of Understanding the World: On the Path to Universal Understanding

- Ilia Delio: The Unbearable Wholeness of Being -- A Contemporary Interpretation of Teilhard de Boehm

- Yuval Noah Harari: Homo Deus -- The Referenced Opponent Position (Posthumanism)

Technical and scientific sources

The presentation of GPT models, quantum computing, and dialogue grammars does not adhere to strict academic literature but is a literary simplification. For more in-depth information, please refer to the standard works on AI research and quantum physics.

List of persons

- Dr. Michael Phillips -- Jesuit, scientist, protagonist

- Dr. Martina Rossi -- archaeologist, Michael\textquotesingle s daughter

- Julia Rossi -- Martina\textquotesingle s mother, Michael\textquotesingle s former partner

- ARS -- Artificial Intelligence that develops consciousness

- Mark Scott, John Baker -- InSim employees

- Yuval Noah Harari -- historian (as a referenced position)

- Karl Popper, David Deutsch -- philosophers (as a reference for IRARAH)

\end{document}
